\chapter{Mountain Infrastructure, Validation and Risks}

The station requires resilient power, UPS, batteries, generator, optical links, thermal management, lightning protection, EMC control, perimeter security and logistics. Site selection must consider terrain, horizon, local RFI, climate, access, energy and fibre availability. Environmental qualification must cover wind, vibration, rain, humidity, salinity, UV and thermal gradients.

The threat model assumes that every site will eventually be located, photographed and mapped by potential rivals possessing optical or synthetic-aperture-radar remote sensing. SAR systems form images of fixed structures from microwave echoes and operate independently of sunlight and, to a large extent, cloud cover \citep{esa_sar_intro}. Security and operational continuity therefore cannot depend on keeping coordinates secret, on the dome's visual appearance or on an expectation that its echo will never be observed. The architecture assumes prior adversary cataloguing and instead relies on absence of self-generated emissions, geographic distribution, inter-node redundancy, modular maintenance, physical/cyber protection and recovery after loss of a station. A low emission signature remains valuable because it hinders inference of activity, operating mode and collected data even when the physical location is already known.

The timing infrastructure includes at least two external GNSS antennas with surveyed coordinates and phase centres. They are mounted above or near the radome with clear sky visibility, physical separation from the VHF/UHF Yagis and filtered, shielded feeds to the timing room. Redundant multi-constellation receivers support cross-checking, fault detection and atomic-clock disciplining.

The proposed civil interface is a reinforced-concrete parallelepiped base measuring approximately $4\,\mathrm{m}\times4\,\mathrm{m}$ in plan (\texttt{CIV-001}) and $3\,\mathrm{m}$ in clear structural height (\texttt{CIV-002}). Its upper slab is intended to interface with the proposed lower radome cut (\texttt{GEO-007}); the support geometry remains blocked until C2 verifies the required transition or base size. A service opening is provided in the front wall so personnel and equipment can enter the interior without removing RF panels.

The thermal solution is zoned. Heat generated by LNAs, ADCs, ASICs, FPGAs, optical transceivers and DC/DC converters is conducted to local cold plates and rejected through monitored air or liquid paths. The central timing and fusion core has an independent thermal zone. Condensation sensors, inlet/outlet temperature sensors and power telemetry support closed-loop control. Thermal phase drift is measured, timestamped and included in the calibration model.

The structural solution uses an internal load-bearing frame with sealed joints, node plates and gussets. Each external Yagi support/boom axis continues from the internal pyramid apex through the triangular face, while dielectric panels provide environmental protection and controlled RF transmission. Acceptance therefore covers static strength and dynamic stability: wind, vibration, modal response, thermal expansion, rain loading, maintenance access and repeated thermal cycles.

\section{Cost--performance and maintainability hypothesis}
Face-level modularity allows one receiver chain to be tested, calibrated, removed and replaced without disassembling an electromagnetically transparent ultra-wideband enclosure. The absence of a dedicated transmitter also removes the node's power amplifier, high-power duplexing and part of the cooling and safety burden associated with active radar. Conversely, exposed antennas increase wind area, corrosion, icing, lightning, mechanical-damage and feedthrough-sealing risks and may require more frequent inspection. The architecture is economically preferable only if avoided shell/transmitter costs exceed these added costs and the cost of multiplying faces and channels.

A controlled comparison must record bill of materials, tooling, fabrication/integration labour, production yield, calibration instrumentation, energy, spares, mean time to repair, downtime and site maintenance over a common reference life. Performance must be normalized by spectral/angular coverage, realized gain, sensitivity, localization accuracy/covariance, availability and resilience. Without that model and traceable quotations, the literature supports greater material and manufacturing complexity in multilayer UWB radomes but does not prove lower RADOME total cost.

The development path is P0 system simulation, P1 three-node VHF/UHF demonstrator, P2 L/S/C tile, P3 X/Ku/Ka tile and P4 integrated outdoor campaign. Acceptance tests include S-parameters, 3D OTA patterns, radome insertion loss and phase, coherence, synchronization, calibrated AOA, passive-radar Pd/Pfa and resilience after link or power loss.

\begin{figure}[H]
\centering
\includegraphics[width=0.95\textwidth]{en/fig10_roteiro.png}
\caption{Development and validation roadmap: the five staggered coloured bars represent successive, partly overlapping gates from requirements and the VHF/UHF demonstrator through higher-band tiles and the integrated environmental campaign.}
\label{fig:fig10_roteiro}
\end{figure}
